\chapter*{Abstract}
The advent of big data and machine learning (ML) has the potential to enhance the quality and timeliness of care for hospitalized patients. Through analyzing time series data of lab results and vital signs, machine learning systems are capable of delivering predictions of the imminent risks of in-hospital mortality or critical disease development. However, the efficacy of such models is hampered by three predominant challenges: (1) current systems often produce too many false alarms, which compromise the model's reliability and acceptance into clinical practice; (2) most ML models struggle to handle the issue of missing data, which is caused by differences in measurement frequencies across channels or by decisions made by clinicians and (3) ML models require expert annotated supervisory labels to train, that are costly and labor-intensive to acquire. 
In this thesis, we address the challenge of high false alarm rates through a novel objective that controls false alarms via a direct constraint on precision. Next, we present probabilistic models to handle missing features in both regularly and irregularly spaced time-series models. For regularly spaced time-series models, we incorporate a probabilistic approach that handles missing features via exact marginalization, thereby avoiding predictions that depend on specific guesses of the missing value. In irregularly spaced time-series models, we propose a generative model that learns the data-dependent missingness mechanism, ultimately yielding superior extrapolative performance compared to alternatives. Finally, to manage limited labels, we leverage a semi-supervised training objective that outperforms alternatives by simultaneously maximizing generative performance across all features (both labeled and unlabeled), while ensuring high-quality predictions on the labeled data. The solutions proposed here address significant challenges in clinical risk prediction models, paving the way for their integration into routine hospital care.